\section{Predicate-Certified Block Navigation}
\label{sec:index}
\subsection{Two Granularities, One Vector Universe}
Prefix-frontier entry changes how search starts; blocks change what it can
search locally.  Physical level 0 contains exact groups and their vector
graphs.  Physical level 1 adds vector graphs over groups aggregated into
blocks.  Within an authorized block, proximity edges can connect vectors
from different exact groups directly, avoiding a requirement to follow the
fine label relation for every transition.  The base remains necessary for
queries that cannot certify a useful coarse region. Figure~\ref{fig:method-overview}
uses the catalog example to connect entry selection with this change of scale.

We call this a \texttt{1L} index: one overlay above the base, hence two
physical levels.  A \texttt{2L} index has levels 0, 1, and 2.  More generally,
an $m$-overlay index has $m+1$ physical levels.  All levels reference the same
vector identifiers, but their block counts, represented membership, and
vector-edge counts can differ.  The levels aggregate label-space regions over shared vectors; the query
algorithm below determines their participation.

\begin{table}[t]
\centering\small
\caption{Canonical topology notation; $L$ means LNG and $T$ means prefix
Trie. Each upper level is an independently built block overlay.}
\label{tab:level-notation}
\begin{tabular}{llll}
\toprule
Code & Level 0 & Level 1 & Level 2\\
\midrule
\texttt{0L[T]} & Trie groups & absent & absent\\
\texttt{1L[T|L]} & Trie groups & LNG blocks & absent\\
\texttt{2L[T|LT]} & Trie groups & LNG blocks & Trie blocks\\
\bottomrule
\end{tabular}
\end{table}

A configuration is $(m,\mathbf{T},\boldsymbol{\tau},e,r)$: overlay count,
increasing thresholds, topology per physical level, entry provider, and
routing policy.  In particular $\tau_0\in\{L,T\}$ is an independent choice;
both base choices are evaluated.  Historical plot labels
\texttt{1L-LNG} and \texttt{2L-LT} denote the fixed-LNG-base controls
\texttt{1L[L|L]} and \texttt{2L[L|LT]}.  The full notation always lists the
base before the separator.

\begin{figure*}[t]
\centering
\begin{tikzpicture}[font=\footnotesize,>=Latex,
 group/.style={draw,rounded corners,minimum width=1.25cm,minimum height=1.08cm,align=center},
 point/.style={circle,draw,inner sep=2pt,fill=white},
 seed/.style={point,draw=blue!75!black,line width=1pt},
 edge/.style={->,thick},lab/.style={align=center}]
% The same four singleton exact groups appear in panels (a) and (b).
\begin{scope}[xshift=0cm]
\node[lab,font=\small\bfseries] at (1.55,2.6) {(a) LNG: one entry};
\node[group] (lb) at (0,1.35) {$\{b\}$\\\phantom{$x_0$}};
\node[group] (lab) at (1.55,1.35) {$\{a,b\}$\\\phantom{$x_1$}};
\node[group] (labc) at (3.1,1.35) {$\{a,b,c\}$\\\phantom{$x_2$}};
\node[seed] (lx0) at (0,1.1) {$x_0$};
\node[point] (lx1) at (1.55,1.1) {$x_1$};
\node[point] (lx2) at (3.1,1.1) {$x_2$};
\draw[edge] (lx0)--(lx1);\draw[edge] (lx1)--(lx2);
\node[group,fill=gray!12] (lac) at (3.1,-0.3) {$\{a,c\}$\\\phantom{$x_3$}};
\node[point,fill=gray!12] (lx3) at (3.1,-0.55) {$x_3$};
\draw[edge,gray] (lac.north)--(labc.south);
\node[lab] at (1.1,-1.35) {Minimal entry: $\{b\}$\\Cross-group edges follow containment.};
\end{scope}
\begin{scope}[xshift=5.5cm]
\node[lab,font=\small\bfseries] at (1.55,2.6) {(b) Prefix: two entries};
\node[group] at (0,1.35) {$\{b\}$\\\phantom{$x_0$}};
\node[group] at (1.55,1.35) {$\{a,b\}$\\\phantom{$x_1$}};
\node[group] at (3.1,1.35) {$\{a,b,c\}$\\\phantom{$x_2$}};
\node[seed] (tx0) at (0,1.1) {$x_0$};
\node[seed] (tx1) at (1.55,1.1) {$x_1$};
\node[point] (tx2) at (3.1,1.1) {$x_2$};
\draw[edge] (tx1)--(tx2);
\node[group,fill=gray!12] at (3.1,-0.3) {$\{a,c\}$\\\phantom{$x_3$}};
\node[point,fill=gray!12] at (3.1,-0.55) {$x_3$};
\node[lab] at (1.55,-1.35) {Frontier: $\{b\},\{a,b\}$\\Both eligible prefix branches are seeded.};
\end{scope}
\begin{scope}[xshift=11cm]
\node[lab,font=\small\bfseries] at (1.4,2.6) {(c) An authorized block};
\draw[rounded corners,fill=green!6] (-0.1,0.35) rectangle (3.05,2.05);
\node[lab] at (1.45,1.78) {Level 1: root $R_b=\{a,b\}$};
\node[seed] (bx1) at (0.65,1.05) {$x_1$};
\node[point] (bx2) at (2.25,1.05) {$x_2$};
\draw[<->,thick,green!40!black] (bx1)--(bx2);
\node[lab,font=\scriptsize] at (1.45,0.58) {Local edges cross group boundaries.};
\node[seed] at (0.4,-0.45) {$x_0$};
\node[lab,anchor=west] at (0.8,-0.45) {Level 0 residual};
\node[lab] at (1.45,-1.35) {$Q=\{b\}$: block authorized.\\$Q'=\{b,c\}$: use finer $\{a,b,c\}$.};
\end{scope}
\end{tikzpicture}
\caption{One query, two entry contracts, and one block overlay. For
$Q=\{b\}$, blue-ringed vertices are seeds and the gray group is ineligible.
Panels (a) and (b) use the same singleton groups. With threshold $T_1=1$,
$x_1,x_2$ form the block in (c); its prefix authorizes either direction of
local navigation for $Q$. Query $Q'$ cannot authorize that block. Arrows denote permitted connections; each state uses one physical
level throughout its traversal.}
\label{fig:method-overview}
\end{figure*}

\subsection{Independent Mass Partitions}
Each terminal has weight equal to its group size.  For threshold $T_\ell$,
a bottom-up pass accumulates uncovered point mass.  A non-root node becomes
a block root when its accumulated mass exceeds $T_\ell$; that mass is then
removed from its parent's uncovered total.  A block owns terminals below its
root up to, but excluding, another block root at the same level.

\begin{algorithm}[t]
\caption{One independent block partition}
\label{alg:partition}
\begin{algorithmic}[1]
\Require canonical trie, terminal weights $w(v)$, threshold $T$
\For{nodes $v$ in postorder}
  \State $u(v)\gets w(v)+\sum_{c\in\mathrm{child}(v)}u(c)$
  \If{$v\ne\mathrm{root}$ and $u(v)>T$}
    \State mark $v$ as a block root; $u(v)\gets0$
  \EndIf
\EndFor
\State assign each terminal to its nearest marked ancestor (including itself), if one exists
\State \Return roots and their direct-member lists $P_b$
\end{algorithmic}
\end{algorithm}

Direct-member sets are disjoint within a layer, but may form only a partial
cover: residual groups below no emitted root remain represented by the base.
Threshold $T_\ell$ is an emission threshold, not a maximum block size.  A
large exact group or branching node can yield a block exceeding it.  Separate
partition passes for $T_1<T_2<\cdots$ leave earlier levels unchanged; they
do not imply nested direct-member sets or identical coverage at all levels.

\subsection{Certification and Connectivity}
Block $b$ stores its root prefix $R_b$, direct members $P_b$, a vector entry
point, and adjacency with an explicit owner level.  It is authorized when
$Q\subseteq R_b$.  Since $R_b\subseteq A_i$ for every $i\in P_b$, the
certificate implies $Q\subseteq A_i$.  An unauthorized block may contain
eligible points; inability to certify it says nothing about those points'
individual labels.

Within a block, the proximity graph need not respect exact-group boundaries.
Between blocks at the same level, either a prefix relation or the LNG cover
relation specifies legal transitions to more specific root-label sets.
Only the layer's direct-member points are used to realize its vector edges.
Both relations preserve eligibility.  Prefix edges are sparse, but are not
necessarily a subset of the LNG \emph{cover edges}: an intermediate observed
set on another prefix branch can remove a cover edge that the prefix forest
contains.  Prefix reachability is a subset of containment reachability.
This is the reason to test topology quality rather than infer it from edge
counts.  Algorithm~\ref{alg:trie-topology} is used on groups at level 0 and
on block roots at an overlay level.

\begin{algorithm}[t]
\caption{Terminal-prefix topology at a fixed physical level}
\label{alg:trie-topology}
\begin{algorithmic}[1]
\Require canonical trie with the level's groups or block roots marked
\State push $(\mathrm{root},\bot)$; $E\gets\varnothing$
\While{stack not empty}
  \State pop $(v,a)$
  \If{$v$ is marked}
    \If{$a\ne\bot$} \State add $(a,v)$ to $E$ \EndIf
    \State $a\gets v$
  \EndIf
  \For{each child $c$ of $v$} \State push $(c,a)$ \EndFor
\EndWhile
\State \Return $E$
\end{algorithmic}
\end{algorithm}

\section{Querying Certified Blocks}
\label{sec:query}
\subsection{Seeding and Optional Routing}
Before graph expansion, a query certifies block roots, applies any routing
gate, and attaches physical levels to its seeds. The retained states then
compete in one bounded queue. Authorization decides which regions are safe;
seeding decides which of them are actually represented in that search.

Entry discovery returns actual group identifiers.  The prefix instance uses
Algorithm~\ref{alg:frontier}; LNG-base controls use one of the two
minimal-superset providers.  These contracts are separate from block
authorization and vector entry-point setup.

We evaluate three routing rules. Let $h$ be the highest authorized
physical level at least 2, and let $M_h$ be the sum of its authorized
blocks' direct-member counts.
\begin{itemize}
\item Ungated: admit the multilayer search.
\item DRH-v1: admit it exactly when such an $h$ exists.
\item DRH-v2: additionally require $M_h\ge T_{\mathrm{gate}}$.
\end{itemize}
A rejected query runs the base with the same topology and entry provider.
In the two-overlay experiments, $h=2$ whenever it exists.

For an accepted multilayer query, authorized blocks contribute their stored
entry points tagged with their physical level.  Each group entry receives
the lowest numbered authorized overlay owning that group, or level 0
if no such overlay exists.  The measured initializer prioritizes block seeds when
trimming to capacity, then fills remaining space with group seeds.  Retagging replaces the base seed, so the initial queue determines
which physical representations participate.

\subsection{Level-Local Search}
Search states are $(i,\delta,\ell)$: vector identifier, query distance, and
physical level.  Visited state is indexed by $(i,\ell)$, while the answer
set is deduplicated by vector identifier.  The shared queue has a configured
active capacity $L_s$.  Algorithm~\ref{alg:query} describes the containment
path with root-prefix authorization used in the reported experiments.

\begin{algorithm}[t]
\caption{Level-local containment search}
\label{alg:query}
\begin{algorithmic}[1]
\Require $(q,Q,K)$, base and overlays $H$, provider $e$, route $r$, capacity $L_s$
\If{$H$ has no overlays}
 \State \Return \Call{BaseSearch}{$q,Q,e,L_s$}
\EndIf
\State $B_Q\gets\{b:Q\subseteq R_b\}$
\If{$r$ is gated}
 \State find highest authorized physical level $h\ge2$, if any
 \If{no such $h$, or ($r=\mathrm{DRH\mbox{-}v2}$ and $M_h<T_{\mathrm{gate}}$)}
  \State \Return \Call{BaseSearch}{$q,Q,e,L_s$}
 \EndIf
\EndIf
\State discover group entries with $e$ and assign seed levels
\State create authorized block seeds and group seeds
\State deduplicate by $(i,\ell)$, mark seeds seen, and score them before trimming
\State prioritize block seeds; trim to $L_s$ and initialize the queue
\While{an unexpanded state $(i,\delta,\ell)$ remains}
 \State mark the closest active unexpanded state expanded (it stays in the queue)
 \State read only level-$\ell$ adjacency of $i$
 \For{each permitted edge to $j$ at level $\ell$}
  \If{$(j,\ell)$ is unseen}
   \State mark it seen; offer $(j,\mathrm{dist}(q,x_j),\ell)$ to the queue
  \EndIf
 \EndFor
\EndWhile
\State \Return the closest at most $K$ distinct vector identifiers in the final active queue
\end{algorithmic}
\end{algorithm}

\paragraph{Queue semantics.}
Candidates are ordered by $(\delta,i,\ell)$, breaking distance ties by
vector identifier and then level. Initialization retains the closest block
seeds first and fills remaining capacity with the closest group seeds. All
collected seeds are scored and marked seen before this selection. Expansion
marks a state without removing it from the active queue. An offered state
enters if capacity is available or it precedes the current worst state;
in the latter case it evicts that state, whether expanded or unexpanded.
Seen marks persist after rejection or eviction. Search ends when the active
queue has no unexpanded state; answers are the closest distinct vector IDs
in that final queue. Thus an evicted state contributes neither a future
expansion nor an answer at its evicted level.

\paragraph{A two-overlay execution example.}
Consider singleton groups with vectors $y_0,\ldots,y_8$ and label sets,
in order, $a,ab,abc,abd,ac$, followed by $acd,ace,acf,ad$.
Concatenation denotes a set under $a<b<c<d<e<f$. For thresholds $T_1=1$ and $T_2=4$, the independent
partitions produce the direct memberships in Table~\ref{tab:seed-trace}.
At level 1, the $a$ block excludes the emitted $ab$ and $ac$ descendants;
at level 2, $a$ owns all nine vectors. Query $Q=\{a\}$ certifies all four
blocks. Its prefix-frontier group $a$ contributes $y_0$, retagged to level 1,
which duplicates that block's seed and is removed during seed deduplication.

\begin{table}[t]
\centering\small
\caption{Illustrative ownership and seed selection, with $K=1$ and $L_s=2$.
Stored block entries and query distances are chosen for this example.}
\label{tab:seed-trace}
\begin{tabular}{clll}
\toprule
Level & Root & Direct members & Entry (distance)\\
\midrule
1 & $a$ & $y_0,y_8$ & $y_0\ (3)$\\
1 & $ab$ & $y_1,y_2,y_3$ & $y_1\ (1)$\\
1 & $ac$ & $y_4,y_5,y_6,y_7$ & $y_4\ (2)$\\
2 & $a$ & $y_0,\ldots,y_8$ & $y_1\ (1)$\\
\bottomrule
\end{tabular}
\end{table}

The two retained seeds are $(y_1,1)$ and $(y_1,2)$, distinct level states of
the same vector. The first expansion uses only level-1 edges. If it exposes
$y_2$ at distance $0.5$, the queue retains $(y_2,1)$ and $(y_1,1)$ and evicts
$(y_1,2)$ before that coarse state expands. This trace shows how extra
representations can compete for a fixed budget even with isolated adjacency.
It also shows why coverage of the full frontier does not imply coverage
after initialization and queue eviction.

\paragraph{Proposition 2: predicate safety.}
Assume all group seeds are eligible and every overlay-tagged group seed is
a direct member of an authorized owner block; base cross-group edges preserve strict
containment; block seeds are direct members of authorized blocks; and each
overlay edge is expanded only with an authorized source owner matching
its state level, and stays within its direct-member block or follows a same-level
strict-superset block transition.  Then every vector whose query distance is
computed in this search is eligible.

\emph{Proof.} Eligible group seeds and certified block seeds satisfy $Q$.
An intra-group edge preserves the label set.  A base cross edge moves to a
superset, preserving $Q$.  Inside a certified block every direct member
satisfies $Q$; a cross-block edge moves to a root containing the source root
and hence $Q$.  Induction over generated states proves the claim. \hfill$\square$

Construction and root-prefix authorization establish these invariants
before traversal, allowing distance evaluation without per-neighbor
predicate tests.

\paragraph{Proposition 3: layer isolation and fallback.}
A successor inherits its source state's physical level, and adjacency owned
by other levels is rejected.  Induction therefore gives level-local
expansion with no promotion or fallthrough.  When the router rejects an
overlay before seeding, deterministic execution with identical base entries
and budget reproduces the base-search path.  The rejected query additionally pays for authorization and routing.

\subsection{Query Cost}
Let $D_q$ be distance evaluations, $E_q$ inspected adjacency entries, $A_q$
candidate insertion attempts, and $P_q$ state expansions. Let
$T_{\mathrm{entry}}$ include group discovery and vector-seed setup, excluding
distances already counted in $D_qd$. All scored seeds, including discarded
ones, contribute to $D_q$.  Direct certification
costs $T_{\mathrm{auth}}=O(\sum_b(|Q|+|R_b|))$ plus metadata traversal.
The measured default array queue can scan its active capacity even when an
insertion is rejected, so a safe bound is
\begin{equation}
 T_q=O\bigl(T_{\mathrm{entry}}+T_{\mathrm{auth}}+D_qd+E_q
                  +(A_q+P_q)L_s+L_s\log L_s\bigr).
\end{equation}
This bound exposes why fewer group edges need not lower latency: the
frontier, required budget, distances, and queue operations can all change.
Per-level visited state, seed buffers, and queue backing storage contribute
query workspace beyond the bounded active queue. A vector reached at different
levels can incur repeated distance evaluations.
